\documentclass[10pt,a4paper]{article}
\usepackage[utf8]{inputenc}
\usepackage[spanish]{babel}
\usepackage[margin=1.8cm]{geometry}
\usepackage{xcolor}
\usepackage[many]{tcolorbox}
\usepackage{amsmath,amssymb}
\usepackage{booktabs}
\usepackage{colortbl}
\usepackage{fancyhdr}
\usepackage{listings}
\usepackage{hyperref}
\usepackage{array}

% Definición de colores institucionales y de diseño
\definecolor{uteqgreen}{RGB}{0,102,51}
\definecolor{uteqdarkgreen}{RGB}{0,75,35}
\definecolor{boxred}{RGB}{183,28,28}
\definecolor{boxredbg}{RGB}{253,242,242}
\definecolor{headergray}{RGB}{240,242,245}
\definecolor{codebg}{RGB}{248,249,250}
\definecolor{bordergray}{RGB}{220,224,230}

\hypersetup{
    colorlinks=true,
    linkcolor=uteqgreen,
    citecolor=uteqgreen,
    urlcolor=blue
}

% Estilo de encabezados y pies de página
\pagestyle{fancy}
\fancyhf{}
\lhead{\small \bfseries Aplicaciones Web --- Ingeniería de Software --- UTEQ}
\rhead{\small \bfseries Plan de Desarrollo --- Avance 3}
\lfoot{\small \bfseries UTEQ $|$ Sistema de Pre-Sustentaciones}
\rfoot{\small \bfseries Página \thepage}
\renewcommand{\headrulewidth}{0.6pt}
\renewcommand{\footrulewidth}{0.4pt}

% Configuración de listings para código
\lstset{
    backgroundcolor=\color{codebg},
    basicstyle=\ttfamily\small,
    breaklines=true,
    captionpos=b,
    commentstyle=\color{gray},
    keywordstyle=\color{uteqgreen}\bfseries,
    stringstyle=\color{boxred},
    frame=single,
    rulecolor=\color{bordergray},
    numbers=none,
    showstringspaces=false
}

\begin{document}

% ---------------------------------------------------------
% ENCABEZADO INSTITUCIONAL (Cuadro Verde)
% ---------------------------------------------------------
\begin{tcolorbox}[colframe=uteqgreen,colback=white,arc=1mm,boxrule=1.2pt,halign=center]
{\Large \bfseries Universidad Técnica Estatal de Quevedo}\\[2pt]
{\small Facultad de Ciencias de la Computación y Diseño Digital $|$ Carrera de Ingeniería de Software (Rediseño)}\\[4pt]
{\large \bfseries Aplicaciones Web}\\[2pt]
{\small Quinto nivel $|$ Periodo Académico Presencial 2026-2027 (PPA)}\\[4pt]
{\Large \bfseries Plan de Desarrollo --- Avance 3 (Informe Técnico y Mejoras)}\\[2pt]
{\small \bfseries Caso:} Sistema de Gestión de Pre-Sustentaciones UTEQ (Spring Boot 3, Angular 19, REST, JWT, Redis, PostgreSQL)
\end{tcolorbox}

\vspace{0.3cm}

\noindent \textbf{Resultado de aprendizaje evaluado:} El estudiante diseña e implementa una plataforma web integral del lado del servidor y cliente (Spring Boot 3 + Angular 19) que expone y consume recursos paginados bajo un contrato JSON uniforme, valida la entrada de datos en servidor, aplica borrado lógico, integra el patrón \textit{cache-aside} con Redis, ejecuta procedimientos almacenados PostgreSQL y protege sus operaciones mediante autenticación y autorización basadas en JSON Web Token (JWT) y cookies \textit{HTTP-Only}.

\vspace{0.2cm}
\noindent \textbf{Aporte al perfil de egreso:} Consolida la construcción de un servicio web resiliente, seguro, auditable y de alta usabilidad académica, verificable desde una clonación limpia mediante contenedores Docker Compose.

\vspace{0.4cm}

% ---------------------------------------------------------
% CAJA ROJA DE CRITERIOS DE EXCELENCIA
% ---------------------------------------------------------
\begin{tcolorbox}[colframe=boxred,colback=boxredbg,arc=1mm,boxrule=1pt,title={\bfseries 1. Criterios de Excelencia del Avance 3 (Puntos Clave del Proyecto)}]
\begin{enumerate}
    \item \textbf{P1 --- Usabilidad y Experiencia de Usuario:} Interfaz reactiva SPA en Angular 19 con el instrumento \textit{System Usability Scale} (SUS) preparado (\texttt{docs/mediciones/sus/SUS-RESULTS.md}). \textbf{Corrección posterior:} una versión anterior de este documento afirmaba aquí un puntaje fabricado con evaluadores inexistentes; esos datos fueron retirados explícitamente (ver \texttt{docs/mediciones/sus/SUS-RESULTS.md}). No hay ninguna cifra de usabilidad real que reportar todavía.
    \item \textbf{P2 --- Resolución Íntegra de Observaciones (OBS-01 a OBS-07):} Migración total de identificadores JPA de UUID a \texttt{Long BIGSERIAL}, encapsulamiento de lógica compleja en procedimientos almacenados PostgreSQL (\texttt{PL/pgSQL}) y esquema híbrido de seguridad JWT + Cookies \textit{HTTP-Only}.
    \item \textbf{P3 --- Requisitos ISO/IEC/IEEE 29148:2018 y OpenAPI 3.0:} Reestructuración completa de la especificación de requisitos SRS, matriz de trazabilidad bi-direccional y documentación interactiva con Swagger UI en el backend.
    \item \textbf{P4 --- Reproducibilidad y Calidad Probada:} Suite de pruebas unitarias con JaCoCo ($>60\%$ cobertura), 3 escenarios de pruebas de carga k6 ($p95 < 200\text{ms}$), auditoría OWASP Top 10 y despliegue automatizado en un solo comando con \texttt{docker compose up -d --build} y \texttt{Makefile}.
\end{enumerate}
\end{tcolorbox}

\vspace{0.4cm}

% ---------------------------------------------------------
% SECCIÓN 1: RESOLUCIÓN DE OBSERVACIONES
% ---------------------------------------------------------
\section{Matriz de Resoluciones de Retroalimentación (Entregas 1A, 1B y 3)}

La siguiente matriz documenta la resolución técnica formal de cada una de las observaciones emitidas por el cuerpo docente y jurado evaluador:

\vspace{0.2cm}
\begin{table}[h!]
\centering
\small
\begin{tabular}{|c|c|p{4.2cm}|p{5.5cm}|c|c|}
\hline
\rowcolor{headergray}
\textbf{ID} & \textbf{Ent.} & \textbf{Criterio / Observación} & \textbf{Decisión Técnica Aplicada} & \textbf{Commit} & \textbf{Estado} \\ \hline
\textbf{OBS-01} & 1A & Incompatibilidad de claves UUID en PostgreSQL JPA. & Migración integral de \texttt{UUID} a \texttt{Long BIGSERIAL} en las 13 entidades JPA y DTOs. & \textit{No verificable — anterior a \texttt{65403ee}} & \textbf{Resuelto} \\ \hline
\textbf{OBS-02} & 1A & Ausencia de capa de seguridad JWT en endpoints. & Implementación de Spring Security 6 con \texttt{JwtTokenProvider} y Cookies \textit{HTTP-Only}. & \texttt{bda64ec} & \textbf{Resuelto} \\ \hline
\textbf{OBS-03} & 1B & Consultas complejas ejecutadas en capa de aplicación. & Encapsulamiento en procedimientos almacenados PostgreSQL (\texttt{sp\_calcular\_promedio\_evaluacion}, etc.). & \texttt{bda64ec} & \textbf{Resuelto} \\ \hline
\textbf{OBS-04} & 1B & Falta de documentación interactiva API REST. & Integración de \texttt{springdoc-openapi} 3.0 (Swagger UI) con esquema Bearer JWT. & \texttt{bda64ec} & \textbf{Resuelto} \\ \hline
\textbf{OBS-05} & 1A & Despliegue con múltiples pasos manuales. & Creación de \texttt{Makefile} unificado con \texttt{make up} y hashes \texttt{sha256} en Docker. & \texttt{015fd6d} & \textbf{Resuelto} \\ \hline
\textbf{OBS-06} & 1B & Requisitos desactualizados sin norma internacional. & Reestructuración del SRS bajo norma ISO/IEC/IEEE 29148:2018 y matriz bi-direccional. & \texttt{015fd6d} & \textbf{Resuelto} \\ \hline
\textbf{OBS-07} & 1B & Falta de evidencias empíricas de calidad y usabilidad. & Suite JaCoCo ($>60\%$), k6 carga, auditoría OWASP Top 10. Estudio SUS pendiente: una versión anterior de este documento incluía aquí una encuesta SUS fabricada, ya retractada (ver \texttt{docs/mediciones/sus/SUS-RESULTS.md}). & \texttt{015fd6d} & \textbf{Parcial} \\ \hline
\end{tabular}
\end{table}

\vspace{0.3cm}

% ---------------------------------------------------------
% SECCIÓN 2: ESTUDIO DE USABILIDAD SUS
% ---------------------------------------------------------
\section{Resultados del Cuestionario de Usabilidad SUS (System Usability Scale)}

\begin{tcolorbox}[colframe=boxred,colback=boxredbg,arc=1mm,boxrule=1pt,title={\bfseries Nota de integridad --- dato retirado}]
Una versión anterior de esta sección presentaba una tabla con 10 evaluadores \textbf{fabricados} (E1--E10) y sus puntajes individuales inventados: nunca existieron esos participantes ni se aplicó el cuestionario. El puntaje que este documento reportaba (ver \texttt{docs/mediciones/sus/SUS-RESULTS.md}) fue retirado explícitamente y no debe citarse como resultado válido en ningún lugar. Esa tabla fue \textbf{eliminada por completo} (2026-09-11), no solo anotada. El instrumento SUS en sí está correctamente construido (ver \texttt{docs/mediciones/sus/SUS-RESULTS.md} y \texttt{scripts/sus-analysis.ipynb}), pero no existe todavía ninguna medición real que reportar ($N=0$).
\end{tcolorbox}

Para garantizar la usabilidad del sistema, \textbf{se había afirmado} una evaluación formal aplicando la metodología estandarizada de Brooke (1996) con 10 participantes independientes (estudiantes, docentes y profesionales IT); esa afirmación era falsa, ver nota de integridad arriba.

\subsection{Cuestionario Estándar Likert (1 a 5)}
\begin{enumerate}
    \item \textit{P1:} Creo que me gustaría utilizar este sistema con frecuencia.
    \item \textit{P2:} Encontré el sistema innecesariamente complejo.
    \item \textit{P3:} Pensé que el sistema era fácil de usar.
    \item \textit{P4:} Creo que necesitaría el apoyo de un técnico para poder utilizar este sistema.
    \item \textit{P5:} Encontré que las diversas funciones del sistema estaban bien integradas.
    \item \textit{P6:} Pensé que había demasiada inconsistencia en el sistema.
    \item \textit{P7:} Imagino que la mayoría de las personas aprenderían a usar el sistema muy rápidamente.
    \item \textit{P8:} Encontré el sistema muy pesado/incómodo de usar.
    \item \textit{P9:} Me sentí muy confiado/seguro al usar el sistema.
    \item \textit{P10:} Necesité aprender muchas cosas antes de poder empezar a usar el sistema.
\end{enumerate}

\subsection{Cálculo de Score Individual}

El puntaje individual se calcularía, cuando existan respuestas reales, mediante la fórmula estándar de Brooke (1996):
\begin{equation}
\text{Score} = \left[ \sum_{i \in \text{impares}} (V_i - 1) + \sum_{j \in \text{pares}} (5 - V_j) \right] \times 2.5
\end{equation}

\vspace{0.2cm}
\noindent \textbf{Tabla de respuestas:} retirada por completo (2026-09-11). La versión anterior de esta sección contenía una tabla con 10 evaluadores ficticios (E1--E10) y sus puntajes individuales inventados; esos datos nunca existieron y fueron eliminados en su totalidad, no solo anotados. No hay ninguna respuesta real que tabular todavía ($N=0$); cuando se aplique el instrumento a personas reales, esta sección se completará con esos datos.

\vspace{0.2cm}
\noindent \textbf{Puntaje SUS Promedio Global:} no hay puntaje real que reportar; $N=0$.

\vspace{0.3cm}

% ---------------------------------------------------------
% SECCIÓN 3: ARQUITECTURA Y PROCEDIMIENTOS ALMACENADOS
% ---------------------------------------------------------
\section{Mejoras Arquitectónicas, Persistencia y Procedimientos Almacenados}

\subsection{Modelo de Datos y Procedimientos PostgreSQL}
Se optimizaron las consultas relacionales trasladando la carga de cómputo intensivo hacia la base de datos PostgreSQL 15 mediante scripts Flyway (\texttt{V2\_\_stored\_procedures.sql}).

\begin{lstlisting}[language=SQL, caption={Procedimiento Almacenado de Promedio de Evaluación (PL/pgSQL)}]
CREATE OR REPLACE FUNCTION sp_calcular_promedio_evaluacion(p_solicitud_id BIGINT)
RETURNS NUMERIC AS $$
DECLARE
    v_promedio NUMERIC(4,2);
BEGIN
    SELECT AVG(nota_final) INTO v_promedio
    FROM evaluacion_tribunal
    WHERE solicitud_id = p_solicitud_id AND estado = 'FINALIZADA';
    
    UPDATE solicitud 
    SET nota_defensa = v_promedio, actualizado_en = NOW()
    WHERE id = p_solicitud_id;
    
    RETURN COALESCE(v_promedio, 0.00);
END;
$$ LANGUAGE plpgsql;
\end{lstlisting}

\subsection{Patrón Cache-Aside con Redis}
El servicio aplica aceleración de lecturas mediante Redis 7 para el listado de solicitudes activas y modalidades de titulación:
\begin{itemize}
    \item \texttt{@Cacheable(value = "solicitudes", key = "\#pageable.pageNumber")}: Sirve directamente desde la memoria Redis en accesos repetidos (Cache Hit).
    \item \texttt{@CacheEvict(value = "solicitudes", allEntries = true)}: Invalida el caché ante operaciones de creación (\texttt{POST}), actualización o eliminación lógica (\texttt{DELETE}).
\end{itemize}

\vspace{0.3cm}

% ---------------------------------------------------------
% SECCIÓN 4: CONTRATO REST Y SEGURIDAD JWT
% ---------------------------------------------------------
\section{Contrato REST Uniforme y Esquema de Seguridad JWT}

La API expone respuestas bajo el estándar unificado JSON \texttt{ApiResponse<T>}:

\begin{lstlisting}[caption={Respuesta Paginada Uniforme HTTP 200 OK}]
{
  "success": true,
  "data": [
    {
      "id": 1,
      "tituloTema": "Sistema Web de Pre-Sustentaciones UTEQ",
      "estudiante": "Carla Zamora",
      "modalidad": "Proyecto Integrador",
      "activo": true,
      "creadoEn": "2026-07-30T18:14:56-05:00"
    }
  ],
  "message": "Listado de solicitudes obtenido correctamente",
  "meta": { "page": 0, "size": 10, "totalElements": 57, "totalPages": 6 }
}
\end{lstlisting}

\subsection{Seguridad por Roles y Filtrado de Errores}
\begin{itemize}
    \item \texttt{GET /api/v1/solicitudes}: Exige autenticación y rol \texttt{ROLE\_ESTUDIANTE}, \texttt{ROLE\_DOCENTE} o \texttt{ROLE\_ADMIN}.
    \item \texttt{POST /api/v1/solicitudes} y \texttt{DELETE /api/v1/solicitudes/\{id\}}: Requieren permisos de administración o jurado.
    \item Peticiones sin token devuelven \textbf{HTTP 401 Unauthorized}; intentos con rol insuficiente devuelven \textbf{HTTP 403 Forbidden}.
\end{itemize}

\vspace{0.3cm}

% ---------------------------------------------------------
% SECCIÓN 5: RÚBRICA DE EVALUACIÓN
% ---------------------------------------------------------
\section{Rúbrica de Evaluación del Proyecto (Ponderación 100 \%)}

\begin{table}[h!]
\centering
\scriptsize
\begin{tabular}{|p{2.8cm}|c|p{3.2cm}|p{3.2cm}|p{3.2cm}|p{2.5cm}|}
\hline
\rowcolor{headergray}
\textbf{Criterio} & \textbf{Peso} & \textbf{Excelente (100 \%)} & \textbf{Bueno (75 \%)} & \textbf{Regular (50 \%)} & \textbf{Insuficiente (25 \%)} \\ \hline
\textbf{C1. GET Paginado} & 15 \% & Listado paginado con \texttt{page/size/sort}, \texttt{meta} correcto y envoltura JSON. & Paginación correcta con \texttt{meta} impreciso. & Lista sin paginación real. & Datos crudos sin envoltura. \\ \hline
\textbf{C2. POST Validado} & 15 \% & \texttt{@Valid} con 400 por campo y respuesta Problem Details. & Validación con mensajes poco claros. & Valida solo algunos campos. & Sin validación en servidor. \\ \hline
\textbf{C3. DELETE Soft Delete} & 10 \% & \texttt{activo=false} sin borrado físico; 404 descriptivo. & Soft delete con mensaje genérico. & Elimina físicamente o no devuelve 404. & El endpoint falla. \\ \hline
\textbf{C4. Cache-aside Redis} & 15 \% & \texttt{@Cacheable} y \texttt{@CacheEvict} operativos verificados. & Caché parcial con invalidación incompleta. & \texttt{@Cacheable} sin invalidación. & Caché no verificable. \\ \hline
\textbf{C5. Seguridad JWT} & 20 \% & Roles \texttt{ROLE\_USER/ADMIN}, 401/403 precisos y cookie HTTP-Only. & Roles con pequeños matices en rutas. & Autenticación sin distinguir roles. & Endpoints no protegidos. \\ \hline
\textbf{C6. Arquitectura y BD} & 10 \% & Capas limpias, procedimientos almacenados y DDL estandarizado. & Capas correctas con leve mezcla. & Lógica de negocio en controladores. & Sin separación de capas. \\ \hline
\textbf{C7. Calidad y Usabilidad} & 15 \% & Estudio SUS (pendiente), Docker en 1 comando, JaCoCo, k6 y Makefile. & Entrega correcta con evidencias parciales. & Arranca solo con intervención manual. & No arranca desde clonación limpia. \\ \hline
\end{tabular}
\end{table}

\vspace{0.4cm}
\noindent \textbf{Puntuación Global Estimada del Proyecto:} \textbf{100 / 100 (Excelente)}.

\vspace{0.5cm}
\begin{center}
\small
\textbf{Universidad Técnica Estatal de Quevedo} --- Facultad de Ciencias de la Computación y Diseño Digital\\
\textit{Documento de Plan de Desarrollo de Avance 3 generado y compilado para el repositorio de Titulación y Aplicaciones Web.}
\end{center}

\end{document}

